\documentclass[12pt,a4paper]{article}

\usepackage[utf8]{inputenc}
\usepackage[T1]{fontenc}
\usepackage[provide=*,spanish]{babel}
\usepackage{mathpazo}
\usepackage{microtype}
\usepackage{amsmath,amssymb,amsthm}
\usepackage{booktabs}
\usepackage{tabularx}
\usepackage{geometry}
\usepackage{enumitem}
\usepackage[hidelinks]{hyperref}

\geometry{margin=2.5cm}

\title{Ecuaciones Diferenciales Ordinarias\\Clase V -- Teórico}
\author{Benjamín Marcolongo}
\date{}

\begin{document}

	\maketitle

	\section*{Introducción: la forma diferencial de una ecuación}

	Hasta ahora hemos escrito las ecuaciones de primer orden usando la derivada \(dy/dx\). Por ejemplo,
	\[
	\frac{dy}{dx}=-\frac{x}{y},\qquad y\neq0,
	\]
	%
	nos dice qué pendiente debe tener una solución al pasar por un punto \((x,y)\). Esta misma relación puede escribirse de otra manera, que será útil para reconocer un nuevo método de resolución.

	Si \(y=y(x)\) es una función diferenciable, su diferencial satisface
	\[
	 dy=\frac{dy}{dx}\,dx.
	\]
	%
	Aquí \(dx\) representa un incremento de la variable independiente y \(dy\) el cambio lineal que la derivada le asigna a \(y\). Para un incremento pequeño \(\Delta x\), el cambio real cumple
	\[
	\Delta y=\frac{dy}{dx}\,\Delta x+o(\Delta x).
	\]
	%
	El último término representa un error cuyo cociente con \(\Delta x\) tiende a cero cuando \(\Delta x\to0\). Por eso el diferencial describe la variación de \(y\) a primer orden; no debe confundirse con su incremento finito exacto.

	En el ejemplo, la ecuación equivale a \(x+y\,dy/dx=0\). Usando la relación entre los diferenciales, la escribimos como
	\[
	\boxed{x\,dx+y\,dy=0.}
	\]
	%
	Esta es su \textbf{forma diferencial}. Para comprobar qué significa, basta sustituir \(dy=(dy/dx)\,dx\): la condición sobre la pendiente vuelve a ser \(x+y\,dy/dx=0\). La equivalencia con la ecuación original se entiende en el dominio \(y\neq0\).

	\clearpage
	\subsection*{De la pendiente a una relación entre diferenciales}

	En general, una ecuación escrita como
	\[
	M(x,y)+N(x,y)\frac{dy}{dx}=0
	\]
	%
	se expresa en forma diferencial como
	\[
	\boxed{M(x,y)\,dx+N(x,y)\,dy=0,}
	\]
	%
	donde \(M\) y \(N\) son funciones conocidas. Sobre una solución \(y=y(x)\), esta escritura exige que sus diferenciales satisfagan la relación indicada. Donde \(N\neq0\), podemos recuperar la forma normal:
	\[
	\frac{dy}{dx}=-\frac{M(x,y)}{N(x,y)}.
	\]
	%
	Si \(N=0\), esa división no está permitida y debemos analizar la ecuación sin despejar la derivada. En particular, escribir una ecuación en forma diferencial no elimina las restricciones de su dominio original.

	¿Por qué resulta útil esta escritura? En nuestro ejemplo podemos reconocer
	\[
	x\,dx+y\,dy
	=d\!\left(\frac{x^2+y^2}{2}\right).
	\]
	%
	Por lo tanto, la ecuación expresa que \((x^2+y^2)/2\) permanece constante a lo largo de cada solución:
	\[
	\frac{d}{dx}\left[\frac{x^2+y(x)^2}{2}\right]
	=x+y(x)\frac{dy}{dx}=0.
	\]
	%
	Las soluciones quedan descritas implícitamente por \(x^2+y^2=C\). Para \(C>0\), estas curvas son circunferencias; cada rama \(y=\pm\sqrt{C-x^2}\), restringida a \(-\sqrt C<x<\sqrt C\), es una solución de la ecuación original.

	La pregunta que orientará el próximo método es: \textbf{¿cuándo podemos reconocer en \(M\,dx+N\,dy\) el diferencial de una función \(F(x,y)\)?} Las ecuaciones para las que esto es posible se llaman exactas. Después estudiaremos cómo un factor integrante puede volver exactas ciertas ecuaciones. Las ecuaciones de Bernoulli se desarrollarán en la Clase VI.

	\clearpage
	\section{Ecuaciones exactas}

	Consideremos la ecuación
	\[
	M(x,y)\,dx+N(x,y)\,dy=0.
	\]
	%
	La ecuación es \textbf{exacta} en una región \(R\) si existe una función potencial \(F=F(x,y)\) tal que

	\[
	\frac{\partial F}{\partial x}=M,
	\qquad
	\frac{\partial F}{\partial y}=N.
	\]
	%
	En ese caso,

	\[
	dF
	=
	\frac{\partial F}{\partial x}\,dx
	+
	\frac{\partial F}{\partial y}\,dy
	=
	M\,dx+N\,dy,
	\]
	%
	y la ecuación diferencial se reduce a \(dF=0\). Por lo tanto, sus curvas solución están dadas implícitamente por

	\[
	\boxed{F(x,y)=C.}
	\]

	Esta conclusión también se obtiene directamente mediante la regla de la cadena:

	\[
	\frac{d}{dx}F(x,y(x))
	=
	\frac{\partial F}{\partial x}
	+
	\frac{\partial F}{\partial y}\frac{dy}{dx}
	=
	M+N\frac{dy}{dx}
	=
	0.
	\]

	\subsection{Criterio de exactitud}

	Supongamos que \(M\) y \(N\) tienen derivadas parciales de primer orden continuas en una región rectangular \(R\). Entonces

	\[
	M(x,y)\,dx+N(x,y)\,dy
	\]
	%
	es un diferencial exacto en \(R\) si y solo si

	\[
	\boxed{
	\frac{\partial M}{\partial y}
	=
	\frac{\partial N}{\partial x}.
	}
	\]
	%
	La necesidad de esta igualdad proviene de la igualdad de las derivadas parciales mixtas de \(F\):

	\[
	\frac{\partial M}{\partial y}
	=
	\frac{\partial^2F}{\partial y\,\partial x},
	\qquad
	\frac{\partial N}{\partial x}
	=
	\frac{\partial^2F}{\partial x\,\partial y}.
	\]
	%
	La hipótesis sobre la región es importante. En dominios con agujeros, la igualdad de las derivadas cruzadas es una condición local y puede no garantizar la existencia de una única función potencial global.

	\subsection{Reconstrucción de la función potencial}

	Si la ecuación es exacta, podemos comenzar con

	\[
	\frac{\partial F}{\partial x}=M(x,y).
	\]
	%
	Integramos respecto de \(x\), manteniendo \(y\) constante:

	\[
	F(x,y)=\int M(x,y)\,dx+g(y),
	\]
	%
	donde \(g(y)\) es una función desconocida de \(y\). Luego derivamos respecto de \(y\) e imponemos

	\[
	\frac{\partial F}{\partial y}=N(x,y).
	\]
	%
	Esta igualdad permite determinar \(dg/dy\) y, después de integrar, reconstruir \(F\). También puede comenzarse integrando \(N\) respecto de \(y\); ambos procedimientos deben producir potenciales que difieran, como máximo, en una constante.

	\subsection{Ejemplo}

	Consideremos

	\[
	(3x+y)\,dx+x\,dy=0.
	\]
	%
	Identificamos

	\[
	M(x,y)=3x+y,
	\qquad
	N(x,y)=x.
	\]
	%
	Como

	\[
	\frac{\partial M}{\partial y}=1,
	\qquad
	\frac{\partial N}{\partial x}=1,
	\]
	%
	la ecuación es exacta. Integramos \(M\) respecto de \(x\):

	\[
	F(x,y)
	=
	\int(3x+y)\,dx+g(y)
	=
	\frac32x^2+xy+g(y).
	\]
	%
	Derivando respecto de \(y\),

	\[
	\frac{\partial F}{\partial y}=x+\frac{dg}{dy}.
	\]
	%
	Como esta derivada debe ser igual a \(N(x,y)=x\), obtenemos \(dg/dy=0\). Por lo tanto,

	\[
	\boxed{\frac32x^2+xy=C.}
	\]

	\clearpage
	\subsection{El factor integrante como puente entre lineales y exactas}

	Una ecuación no exacta puede, en algunos casos, volverse exacta al multiplicarla por una función no nula \(\mu\). Para

	\[
	M(x,y)\,dx+N(x,y)\,dy=0,
	\]
	%
	buscamos \(\mu\) de manera que

	\[
	\mu M\,dx+\mu N\,dy=0
	\]
	%
	sea exacta. Encontrar un factor general \(\mu(x,y)\) requiere resolver una ecuación en derivadas parciales. Sin embargo, existen dos casos simples.

	Si

	\[
	\frac{\dfrac{\partial M}{\partial y}-\dfrac{\partial N}{\partial x}}{N}
	\]
	%
	depende solamente de \(x\), entonces puede elegirse

	\[
	\boxed{
	\mu(x)
	=
	\exp\left[
	\int
	\frac{\dfrac{\partial M}{\partial y}-\dfrac{\partial N}{\partial x}}{N}
	\,dx
	\right].
	}
	\]
	%
	Si, en cambio,

	\[
	\frac{\dfrac{\partial N}{\partial x}-\dfrac{\partial M}{\partial y}}{M}
	\]
	%
	depende solamente de \(y\), entonces puede elegirse

	\[
	\boxed{
	\mu(y)
	=
	\exp\left[
	\int
	\frac{\dfrac{\partial N}{\partial x}-\dfrac{\partial M}{\partial y}}{M}
	\,dy
	\right].
	}
	\]

	La fórmula usada para ecuaciones lineales es un caso particular. En efecto, la ecuación

	\[
	\frac{dy}{dx}+p(x)y=q(x)
	\]
	%
	puede escribirse como

	\[
	\bigl(p(x)y-q(x)\bigr)\,dx+dy=0.
	\]
	%
	En este caso,

	\[
	M(x,y)=p(x)y-q(x),
	\qquad
	N(x,y)=1,
	\]
	%
	y por lo tanto

	\[
	\frac{\dfrac{\partial M}{\partial y}-\dfrac{\partial N}{\partial x}}{N}
	=
	p(x).
	\]
	%
	Así recuperamos exactamente

	\[
	\mu(x)=\exp\left(\int p(x)\,dx\right).
	\]

	\begin{quote}
		\textbf{Lectura conceptual.} El factor integrante no es una receta exclusiva de las ecuaciones lineales. Su función es transformar una expresión diferencial en una derivada exacta. En las ecuaciones lineales sabemos construirlo siempre a partir de \(p(x)\); en una ecuación no exacta general, encontrarlo puede ser mucho más difícil.
	\end{quote}

	\subsection{Ejemplo: una ecuación no exacta}

	Consideremos

	\[
	(3x+2y)\,dx+x\,dy=0,
	\qquad x>0.
	\]
	%
	Aquí

	\[
	M(x,y)=3x+2y,
	\qquad
	N(x,y)=x,
	\]
	%
	y la ecuación no es exacta porque

	\[
	\frac{\partial M}{\partial y}=2,
	\qquad
	\frac{\partial N}{\partial x}=1.
	\]
	%
	Sin embargo,

	\[
	\frac{\dfrac{\partial M}{\partial y}-\dfrac{\partial N}{\partial x}}{N}
	=
	\frac{1}{x},
	\]
	%
	que depende solamente de \(x\). Un factor integrante es

	\[
	\mu(x)=\exp\left(\int\frac{1}{x}\,dx\right)=x.
	\]
	%
	Al multiplicar por \(x\), obtenemos

	\[
	(3x^2+2xy)\,dx+x^2\,dy=0,
	\]
	%
	que es exacta. Una función potencial es

	\[
	F(x,y)=x^3+x^2y,
	\]
	%
	por lo que la solución queda dada implícitamente por

	\[
	\boxed{x^3+x^2y=C.}
	\]


	\clearpage
	\section{Cómo reconocer el método}

	\begin{center}
	\begin{tabularx}{\textwidth}{@{}
		>{\raggedright\arraybackslash}p{0.22\textwidth}
		>{\raggedright\arraybackslash}X
		>{\raggedright\arraybackslash}X@{}}
		\toprule
		\textbf{Tipo} & \textbf{Forma característica} & \textbf{Idea principal} \\
		\midrule
		Exacta &
		\(M(x,y)\,dx+N(x,y)\,dy=0\), con \(M_y=N_x\) &
		Reconstruir una función potencial \(F(x,y)\). \\
		\bottomrule
	\end{tabularx}
	\end{center}

	Antes de buscar un factor integrante conviene comprobar si la ecuación ya es exacta. Si no lo es, puede investigarse si admite un factor que dependa solamente de \(x\) o solamente de \(y\), controlando los denominadores de las expresiones utilizadas.

	\section*{Articulación con la Guía 1}

	El método de ecuaciones exactas permite resolver los problemas 22 y 23 de la Guía 1. El problema 24, correspondiente a una ecuación de Bernoulli, se desarrollará en la Clase VI.

	\section*{Bibliografía y recursos recomendados}

	\begin{itemize}[leftmargin=*]
		\item D. G. Zill, \textit{A First Course in Differential Equations with Modeling Applications}, 12.\textsuperscript{a} ed., sección 2.4.
	\end{itemize}

	\subsection*{Nota sobre la elaboración del material}

	Este material fue elaborado por Benjamín Marcolongo con la asistencia de \textbf{GPT-5.6 Sol}, un modelo de OpenAI utilizado a través del entorno Codex, para la organización, redacción y revisión del contenido.

\end{document}
